\section{Setup time can break minimum-requirement replacement}
\label{app:setupmatchingboundary}
The setup-free theorem replaces every same-cardinality matching by one
with minimum summed early requirement. Positive setup can invalidate
that exact replacement even when the assignment graph is complete and
both perfect matchings remain feasible. Here setup is charged for each
nonzero early charge/export action separately on each bus and occupies the shared
connector; it is not merely a penalty for changing charge/export direction.

\begin{proposition}[A setup-dependent assignment boundary]
\label{prop:setupmatchingboundary}
Take $n=2$, $B=\alpha=\delta=1$, $a=b=(9/10,2/5)$,
the complete successor graph, early directional powers $K_c=K_d=1$,
terminal power/budget $H=2$, fee $f=1$ and setup time $\tau=3/20$.
Retain the common windows, full inventory, zero other consumption and
no terminal setup. The uniquely minimum-requirement perfect matching
cannot reproduce every feasible load--cost point of the other matching.
\end{proposition}

\begin{proof}
The two perfect matchings have the following requirements:
\begin{center}
\begin{tabular}{lcc}
\toprule
Matching & Summed early requirement & Mandatory charging recipients\\
\midrule
Diagonal & $4/5$ & One\\
Cross & $3/5$ & Two, each requiring $3/10$\\
\bottomrule
\end{tabular}
\end{center}
On the diagonal, purchase $C=4/5$, $X=0$ on the high bus in one
early action. Setup plus charge takes $3/20+4/5=19/20$.
After A its SOC is $1/10$; charging raises it to $9/10$, and its
B service leaves zero. The low bus moves from $3/5$ after A to $1/5$
after B without early charging. Terminal purchases $1$ and $4/5$
restore full inventory, using $(9/5)/2=9/10$ of the terminal window.
Thus $(E,L,c)=(4/5,9/5,2)$ is physically feasible.

The cross matching requires positive early charge on two different buses.
At $C=4/5,X=0$ it needs at least $4/5+2(3/20)=11/10$ early time.
Splitting actions or lowering power cannot help. Nor can cycling rescue
the net-load point: $E=4/5$ implies $C=4/5+X$, increasing gross use
when $X>0$, while both mandatory recipients remain. Setup cannot overlap
another action on the shared connector, and terminal charging cannot
repair an earlier service inventory violation.

The cross matching is nevertheless feasible at $C=3/5,X=0$: charge
$3/10$ on each bus in total early time $9/10$, then purchase $2$ terminal
units at power $2$. A partial matching or paid idle bus cannot reproduce
$c=2$, since the two simultaneous services in each batch force at least
two buses, and any such extra-bus alternative costs at least $3$.
\end{proof}

A monotone linear query also exposes the omission. For $F(E,L)=E+2L$,
full unit-efficiency refill gives $E+L=13/5$, so total objective is
$c+26/5-E$. Any early action leaves at most $17/20$ energy-transfer
time; the diagonal attains $E=17/20$ with one high-bus charge.
The cross requires two actions and attains at most $E=7/10$, realized
by purchases $3/10$ and $2/5$. Their optimal values are therefore
$127/20$ and $13/2$. Every fleet costing at least $3$ has value at least
$147/20$, so lower-cardinality alternatives do not repair the restricted
minimum. The error from retaining only minimum-$r_m$ matchings is $3/20$.
Because the objective is linear, the true physical and hull values agree;
this is not a positive support-gap example.

The setup-free outer polygons and the slack-anchor inner bounds in
Appendix~\ref{app:setuptime} remain valid. What fails is exact matching
replacement: concentration of recharge among buses matters when each
recipient/action consumes setup time. This supplies neither a minimal
replacement statistic nor a hardness, public-timetable, empirical or
learning claim.
For charge-only operation, Appendix~\ref{app:setupchargeonly} gives exact
classical fixed-matching and joint-duty baselines that retain this setup
information explicitly.
